\section{Risk certificates for the constructive estimator}

The risk analysis separates estimation of cell membership probabilities from estimation of arrived-cell outcome means. Independence holds in the uncapped Poisson comparison experiment used for this analysis. The reported finite-sample estimator instead averages assignments of a capped prefix, and the prefix-transfer argument below connects its risk to that independent-stream comparison. The polynomial branch controls the contribution of sparsely observed cells. The resulting certificate bounds squared error for the observed-cell functional in \cref{obj:def:cell-functional}; identification then connects that certificate to the causal upper bound in \cref{obj:thm:unrestricted-mixed-count-upper}.

For the deterministic mixed-count estimator \leanref{sym:\widehat\tau^{\mathrm{mix}}_{n,d,q}}{\ensuremath{\widehat\tau^{\mathrm{mix}}_{n,d,q}}} of \cref{obj:def:mixed-count-estimator}, write \(N=nq\) for the effective sample size, \(m=n/6\) for the intensity of each auxiliary stream, and \(N_0=mq=N/6\) for the effective stream size. The logarithmic scale \(\ell\), degree \(k\), and threshold \(B\) are specified in \cref{obj:synth_7,obj:synth_8,obj:synth_18}. Define the prefix-tail constant by
\[
\gamma_0=\log 2-\frac12.
\]
On the polynomial branch, where \(N>1\), \(\ell\ge128\), and \(d\ge\sqrt N\), introduce the certificate quantities
\[
b_{n,d,q}
=2\left\{\frac{4dB}{N_0k^2}+3e^{-16\ell}\right\}
\]
and
\[
v_{n,d,q}
=8\left(\frac4{N_0}+\frac1m\right)
+64d(e^{\ell/8}+1)
\left(\frac{B^2}{N_0^2}+\frac{B}{mN_0}\right)
+32e^{-32\ell}\left(1+\frac1m\right).
\]
Here \(b_{n,d,q}\) is the bias allowance and \(v_{n,d,q}\) is the additive second-moment allowance. These quantities summarize the polynomial approximation, marked-count fluctuations, and estimation of membership weights.

The deterministic squared-error envelope \leanref{sym:\mathfrak U_{n,d,q}}{\ensuremath{\mathfrak U_{n,d,q}}} combines these allowances with the elementary branch and the prefix-overflow contribution. The following statement gives its full formula and compares it with the rate \leanref{sym:r_{n,d,q}}{\ensuremath{r_{n,d,q}}} of \cref{obj:def:frontier-rate} throughout the unrestricted-arrival class \leanref{sym:\mathcal M^{+}_{n,d,q}}{\ensuremath{\mathcal M^{+}_{n,d,q}}} of \cref{obj:def:unrestricted-arrival-model-class}.

\begin{theoremv}{thm:unrestricted-uniform-mixed-count-certificate}[Uniform mixed-count risk certificate]
There exists a universal constant \(0<\overline C<\infty\) such that the following holds.
\begin{itemize}
\item \textbf{(Domain.)} Let \(n,d\ge1\) be integers and let \(0<q\le1\).
\item \textbf{(Model.)} Let \(P\in\mathcal M^{+}_{n,d,q}\), the unrestricted-arrival model class of \cref{obj:def:unrestricted-arrival-model-class}.
\item \textbf{(Sampling.)} Let \(\mathbf O_n\), the sample of observed records, satisfy \cref{obj:ass:iid-sampling} under \(P\).
\end{itemize}
Use the mixed-count estimator of \cref{obj:def:mixed-count-estimator}. Let \(\ell\), the logarithmic scale, be defined by \cref{obj:synth_7}; let \(k\), the polynomial degree, be defined by \cref{obj:synth_8}; and let \(B\), the threshold, be defined by \cref{obj:synth_18}. Put
\[
N=nq,\qquad m=\frac n6,\qquad N_0=mq=\frac N6,
\qquad \gamma_0=\log 2-\frac12.
\]
Define the deterministic envelope \(\mathfrak U_{n,d,q}\) as follows. If \(N\le1\), set
\[
\mathfrak U_{n,d,q}=1.
\]
If \(N>1\) and either \(\ell<128\) or \(d<\sqrt N\), set
\[
\mathfrak U_{n,d,q}
=\min\left\{4,
\left(\frac{8d}{eN_0}\right)^2
+4\left(\frac4{N_0}+\frac1m\right)
+4e^{-n\gamma_0}\right\}.
\]
If \(N>1\), \(\ell\ge128\), and \(d\ge\sqrt N\), put
\[
b_{n,d,q}
=2\left\{\frac{4dB}{N_0k^2}+3e^{-16\ell}\right\},
\]
\[
v_{n,d,q}
=8\left(\frac4{N_0}+\frac1m\right)
+64d(e^{\ell/8}+1)
\left(\frac{B^2}{N_0^2}+\frac{B}{mN_0}\right)
+32e^{-32\ell}\left(1+\frac1m\right),
\]
and set
\[
\mathfrak U_{n,d,q}
=\min\left\{4,b_{n,d,q}^2+v_{n,d,q}
+4e^{-n\gamma_0}\right\}.
\]
Then, with \(\Phi(P)\) the cell functional of \cref{obj:def:cell-functional} and \(r_{n,d,q}\) the rate of \cref{obj:def:frontier-rate}, expectation under the canonical i.i.d. observed-sample law satisfies
\[
E_P\!\left[
\left(\widehat\tau^{\mathrm{mix}}_{n,d,q}(\mathbf O_n)-\Phi(P)\right)^2
\right]
\le \mathfrak U_{n,d,q}
\le \overline C\,r_{n,d,q}.
\]
\end{theoremv}

The certificate distinguishes information from arrived outcomes, measured by \(N_0\), from information about cell membership, measured by \(m\). On the polynomial branch, the squared bias allowance captures the aggregate effect of estimating means in sparsely observed cells, while the additive allowance accounts for fluctuations in both components of the cell contribution. The elementary branch supplies the corresponding bound in its stated parameter range. Together, the branches give a deterministic precision benchmark depending on \(n,d,q\), with a universal comparison constant.

\subsection{Polynomial bias and marked-count second moments}

The polynomial statistic \(H_j\) and the pilot-selected statistic \(G_j\) are specified in \cref{obj:synth_14,obj:synth_15}. Their mathematical role is to replace the elementary cell statistic where the pilot count indicates sparse outcome arrival. The degree and threshold jointly determine the approximation contribution \(4dB/(N_0k^2)\) in the bias allowance. The exponentially small term in that allowance accounts for branch control. The certificate combines these contributions before squaring, preserving the distinction between bias and second-moment terms.

Marked counts retain the association between an arrived outcome and its cell. Consequently, the second-moment analysis concerns the marked polynomial statistic, together with the pilot-selected branch, rather than an unmarked count alone. The expression \(v_{n,d,q}\) records the resulting allowance: its terms involve the effective arrived-outcome intensity, the membership-stream intensity, the alphabet size, and the polynomial tuning parameters. This decomposition explains why the degree enters the bias allowance while the coefficient growth associated with the polynomial construction enters the additive allowance.

\subsection{Membership weights and branch control}

The membership weight \(W_j\) in \cref{obj:synth_11} is formed from the membership-labelled records. In the uncapped Poisson comparison experiment, that stream is independent of the pilot and estimation streams. The finite-sample estimator uses the capped prefix and conditional averaging in \cref{obj:def:mixed-count-estimator}; the overflow allowance and prefix-transfer comparison connect it to the independent-stream analysis. The weight supplies the population mass attached to each arrived-cell mean while retaining information from all observed memberships.

The two sample scales appear explicitly in the certificate through \(1/N_0\), \(1/m\), and the mixed term \(B/(mN_0)\). These terms retain the cost of estimating membership weights alongside the cost of estimating arrived-cell means. The pilot-selected statistic in \cref{obj:synth_15} organizes the latter task within the polynomial branch, whereas the elementary ratio in \cref{obj:synth_12} supplies the estimator on the elementary branch. The branch conditions in \cref{obj:thm:unrestricted-uniform-mixed-count-certificate} specify where each deterministic expression applies.

\subsection{Transfer to the fixed observed sample and causal target}

The auxiliary construction in \cref{obj:def:mixed-count-estimator} uses an independent Poisson prefix and uniform stream assignments. It assigns zero output when the prefix exceeds the available \(n\) records and clips the output to \([-1,1]\). The term \(4e^{-n\gamma_0}\) in \cref{obj:thm:unrestricted-uniform-mixed-count-certificate} is the prefix-overflow allowance. Its dependence on the full record count \(n\) reflects the supply of records for the auxiliary construction, while \(N=nq\) governs the effective outcome information. Clipping supplies the bounded-loss allowance represented by the minimum with \(4\); the zero estimator on the small-effective-sample branch supplies the bound \(1\).

Conditional averaging over the auxiliary randomization produces the deterministic estimator in \cref{obj:def:mixed-count-estimator}. The squared-loss comparison for such averaging is supplied by \cref{obj:lem:randomized-kernel-barycenter}. Thus the certificate concerns the estimator evaluated on the fixed observed sample, with the auxiliary construction serving its risk analysis.

For the causal interpretation, \cref{obj:lem:unrestricted-cell-identification} identifies the observed-cell functional with the average treatment effect under the unrestricted model conditions. Substituting that identity into \cref{obj:thm:unrestricted-uniform-mixed-count-certificate} gives the causal risk control used in \cref{obj:thm:unrestricted-mixed-count-upper}. The same identification step applies throughout the rare-arrival class \leanref{sym:\mathcal M_{n,d,q}}{\ensuremath{\mathcal M_{n,d,q}}} of \cref{obj:def:model-class}. The following specialization records the resulting envelope for that class.

\begin{theoremv}{thm:uniform-mixed-count-certificate}[Uniform mixed-count risk certificate]
There exists a universal constant \(0<\overline C<\infty\) such that the following holds for every \(n,d\in\mathbb N\), \(q\in\mathbb R\), and full-data law \(P\in\mathcal M_{n,d,q}\) satisfying \cref{obj:def:model-class}, including \(n\ge1\), \(d\ge1\), and \(0<q\le1\).

Let \(\widehat\tau^{\mathrm{mix}}_{n,d,q}\) be the estimator in \cref{obj:def:mixed-count-estimator}, and let \(r_{n,d,q}\) be the rate in \cref{obj:def:frontier-rate}. Let \(\ell\) be the logarithmic scale defined in \cref{obj:synth_7}, \(k\) the polynomial degree defined in \cref{obj:synth_8}, and \(B\) the threshold defined in \cref{obj:synth_18}. Put
\[
N=nq,\qquad m=\frac n6,\qquad N_0=mq=\frac N6,\qquad
\gamma_0=\log 2-\frac12,
\]
where \(N\) is the effective sample size, \(m\) the stream intensity, \(N_0\) the effective stream size, and \(\gamma_0\) the prefix-tail constant. Define the deterministic envelope \(\mathfrak U_{n,d,q}\) by the following branches:
\begin{itemize}
\item \textbf{(Small effective sample.)} If \(N\le1\), set
\[
\mathfrak U_{n,d,q}=1.
\]
\item \textbf{(Elementary branch.)} If \(N>1\) and either \(\ell<128\) or \(d<\sqrt N\), set
\[
\mathfrak U_{n,d,q}
=\min\left\{4,
\left(\frac{8d}{eN_0}\right)^2
+4\left(\frac4{N_0}+\frac1m\right)
+4e^{-n\gamma_0}\right\}.
\]
\item \textbf{(Polynomial branch.)} If \(N>1\), \(\ell\ge128\), and \(d\ge\sqrt N\), define the certificate quantities
\[
b_{n,d,q}
=2\left(\frac{4dB}{N_0k^2}+3e^{-16\ell}\right),
\]
\[
v_{n,d,q}
=8\left(\frac4{N_0}+\frac1m\right)
+64d\left(e^{\ell/8}+1\right)
\left(\frac{B^2}{N_0^2}+\frac{B}{mN_0}\right)
+32e^{-32\ell}\left(1+\frac1m\right),
\]
and set
\[
\mathfrak U_{n,d,q}
=\min\left\{4,b_{n,d,q}^2+v_{n,d,q}+4e^{-n\gamma_0}\right\}.
\]
\end{itemize}
Then, writing \(\mathbf O_n\) for the observed sample induced by \(P\) and \(\tau(P)\) for the average treatment effect in \cref{obj:def:ate-functional},
\[
E_P\!\left[
\left(\widehat\tau^{\mathrm{mix}}_{n,d,q}(\mathbf O_n)-\tau(P)\right)^2
\right]
\le \mathfrak U_{n,d,q}
\le \overline C\,r_{n,d,q}.
\]
\end{theoremv}

The rare-arrival restriction places the model within the domain of the unrestricted certificate, and \cref{obj:lem:cell-identification-background} fixes the causal target there. Accordingly, the specialization retains the same numerical envelope. Its polynomial bias, second-moment, and prefix-transfer components are those of \cref{obj:thm:unrestricted-uniform-mixed-count-certificate}; model inclusion and identification supply the specialization step.

The next statement records this causal upper bound together with the squared-loss comparison between the deterministic estimator and its auxiliary randomized output. It makes explicit the role of conditional averaging in the final estimator.

\begin{theoremv}{thm:mixed-count-upper}[Mixed-count risk upper bound]
There exists a universal constant \(0<\overline C<\infty\) such that the following holds for every choice of parameters and law satisfying:
\begin{itemize}
\item \textbf{(Sample size and dimension.)} \(n,d\in\mathbb N\) satisfy \(n,d\ge1\).
\item \textbf{(Arrival floor.)} \(q\in\mathbb R\) satisfies \(0<q\le1\).
\item \textbf{(Model.)} The full-data law \(P\) belongs to \(\mathcal M_{n,d,q}\) in \cref{obj:def:model-class}.
\end{itemize}
Let \(\widehat\tau^{\mathrm{mix}}_{n,d,q}\) and \(\widetilde\tau(\mathbf O_n;\omega)\) be the deterministic estimator and auxiliary randomized output of \cref{obj:def:mixed-count-estimator}, where \(\mathbf O_n\) is the observed sample and \(\omega\) is the auxiliary randomization. Let \(\ell\) be the logarithmic scale defined in \cref{obj:synth_7}, \(k\) the polynomial degree defined in \cref{obj:synth_8}, and \(B=256\ell\) the threshold defined in \cref{obj:synth_18}. Use the elementary and polynomial branch rules of \cref{obj:def:mixed-count-estimator}. Define the effective sample size \(N\), stream intensity \(m\), effective stream size \(N_0\), and prefix-tail constant \(\gamma_0\) by
\[
N=nq,\qquad m=\frac n6,\qquad N_0=mq=\frac N6,
\qquad \gamma_0=\log 2-\frac12.
\]
Define the certificate quantities
\[
b_{n,d,q}
=
2\left(\frac{4dB}{N_0k^2}+3e^{-16\ell}\right)
\]
and
\[
v_{n,d,q}
=
8\left(\frac4{N_0}+\frac1m\right)
+
64d\left(e^{\ell/8}+1\right)
\left(\frac{B^2}{N_0^2}+\frac{B}{mN_0}\right)
+
32e^{-32\ell}\left(1+\frac1m\right).
\]
The deterministic risk envelope \(\mathfrak U_{n,d,q}\) is defined as follows. For \(N\le1\), set
\[
\mathfrak U_{n,d,q}=1.
\]
For \(N>1\) on the elementary branch, set
\[
\mathfrak U_{n,d,q}
=
\min\left\{
4,\,
\left(\frac{8d}{eN_0}\right)^2
+4\left(\frac4{N_0}+\frac1m\right)
+4e^{-n\gamma_0}
\right\}.
\]
For \(N>1\) on the polynomial branch, set
\[
\mathfrak U_{n,d,q}
=
\min\left\{
4,\,
b_{n,d,q}^2+v_{n,d,q}+4e^{-n\gamma_0}
\right\}.
\]
Then, with \(\tau(P)\), the average treatment effect, defined in \cref{obj:def:ate-functional}, and \(r_{n,d,q}\), the frontier rate, defined in \cref{obj:def:frontier-rate},
\[
E_P\!\left[
\left(
\widehat\tau^{\mathrm{mix}}_{n,d,q}(\mathbf O_n)-\tau(P)
\right)^2
\right]
\le
\mathfrak U_{n,d,q}
\le
\overline C\,r_{n,d,q}.
\]
The deterministic estimator also satisfies
\[
E_P\!\left[
\left(
\widehat\tau^{\mathrm{mix}}_{n,d,q}(\mathbf O_n)-\tau(P)
\right)^2
\right]
\le
E_{P,\omega}\!\left[
\left(
\widetilde\tau(\mathbf O_n;\omega)-\tau(P)
\right)^2
\right],
\]
where the auxiliary expectation uses the independent Poisson prefix and uniform stream assignments, including the zero output when the prefix exceeds \(n\), specified in \cref{obj:def:mixed-count-estimator}.
\end{theoremv}

Averaging over the prefix and stream assignments removes auxiliary variation under squared loss. The first inequality in \cref{obj:thm:mixed-count-upper} is the certificate of \cref{obj:thm:uniform-mixed-count-certificate}, and the final comparison is the conditional-averaging relation of \cref{obj:lem:randomized-kernel-barycenter} applied to the specified auxiliary output. These references preserve the same construction and envelope across the two statements. For a prospective choice of \(n,d,q\), the envelope therefore supplies a deterministic causal risk benchmark that incorporates cell complexity, outcome arrival, and the fixed-sample implementation.
